\subsection{Optionen und Aktionen}\label{subsec:optionen-und-aktionen}

Im folgenden Abschnitt werden die Aktionen und Optionen, die der Spieler im Spiel verwenden kann, aufgelistet und beschrieben.\\



\let\origclearpage\clearpage
\renewcommand{\clearpage}{\relax}

\newgeometry{top=2.5cm, bottom=2.5cm, left=1.5cm, right=1.5cm}

\begin{longtable}{|p{1.6cm}|p{1.9cm}|p{5.5cm}|p{3.2cm}|p{4.2cm}|}
    \hline
    \rowcolor{gray!30}
    \textbf{ID / Name} & \textbf{Akteure} & \textbf{Ereigniseinfluss} & \textbf{Anfangsbedingungen} & \textbf{Abschlussbedingun- gen} \\
    \hline
    \endfirsthead

    \hline
    \rowcolor{gray!30}
    \textbf{ID / Name} & \textbf{Akteure} & \textbf{Ereigniseinfluss} & \textbf{Anfangsbedingungen} & \textbf{Abschlussbedingun- gen} \\
    \hline
    \endhead

    % Inhalt der Tabelle
    ID01: Laufen & Spieler, Polizei KI, NPC KI&
    1. Spieler drückt Taste (oben, unten, rechts, links / w, a, s, d) \newline
    2. Spielfigur bewegt sich in die entsprechende Richtung &
    Spielfigur befindet sich auf der Map &
    Spielfigur befindet sich weiter rechts/links/ oben/unten auf der Map \\
    \hline
    ID02: mit Tür interagieren & Spieler &
    1. beide Spielfiguren befinden sich in unmittelbarer Nähe einer begehbaren Tür \newline
    2. ein Spieler drückt die Interagieren-Taste\newline
    3. Tür öffnet sich und ein weiterer Teil der Map ist begehbar &
    beide Spielfiguren befinden sich in der Nähe einer geschlossenen Tür &
    die Spielfiguren befinden sich in der Nähe einer offenen Tür\\
    \hline
    ID03: schießen & Spieler, Kartell KI &
    1. Spieler drückt Taste \newline
    2. Spielfigur schiesst in die die Richtung in die sie schaut \newline
    3. Jeder NPC (ausser Verkäufer) und Mülleimer, der im Sichtkegel ist erhält schaden (verringert HP) &
    Spielfigur muss eine Waffe ausgewählt haben &
    Jede Spielfigur, die im Sichtkegel ist hat schaden erhalten (verringerte HP)
    ODER Mülleimer fällt um und dropt manchmal ein Item \\
    \hline
    ID04: schlagen & Spieler, Kartell KI &
    1. Spieler drückt Taste \newline
    2. Spielfigur führt Schlag aus \newline
    3. Mülleimer/Gegner in unmittelbarem Umfeld kann Schaden nehmen \newline
    Gegner verliert HP \newline
    ODER\newline
    Mülleimer fällt um und dropt manchmal ein Item&
    Spielfigur hat keine Waffe ausgewählt &
    Mülleimer ist liegt auf dem Boden ggf. ein Item daneben \newline
    ODER\newline
    Gegner/Spielfigur hat weniger HP \\
    \hline
    ID05: kochen & Spieler &
    1. Spieler drückt die interagieren Taste, während beide Spielfiguren sich in der Nähe einer Crafting Station befindet \newline
    2. Spieler wählt Rezep aus, welches zubereitet werden soll.
     Ein Minispiel startet, was erfolgreich beendet werden muss um die Drogen herzustellen.
     In diesem werden zufällig verschiedene Buttons in verschiedenen Farben auf dem Bildschirm angezeigt, 
     die der Spieler mit der zugehörigen Farbe dann in einer gewissen Zeit drücken muss. 
     Wenn die Spieler zu oft zu langsam sind oder die falsche Taste drücken schlägt die Herstellung fehl 
     und das Minispiel endet und die zur Herstellung verwendeten Materialien werden zerstört. &

    Beide Spielfiguren befinden sich in der Nähe der Crafting Station &
    Weniger Ressourcen und falls das Kochen erfolgreich war mehr Drogen im Inventar \\
    \hline
    ID06: mit Käufer/ Verkäufer handeln & Spieler &
    1. Spieler drückt interagieren Taste, während die Spielfigur sich in der Nähe eines interaktiven NPC aufhält \newline
    2. je nach NPC: \newline
    - Verkäufer: es öffnet sich ein Dialog-Feld und zeigt die zu kaufenden Gegenstände an \newline
    - Käufer: bietet bestimmte Menge Geld für bestimmte Menge Drogen (unterschiedlich je Käufer), Spieler kann das Angbot annehmen oder ablehnen &
    Spielfigur befindet sich in unmittelbarer Nähe von einem interaktiven NPC &
    Bei Kauf: Mehr Ressourcen und weniger Geld \newline
    ODER \newline
    Bei Verkauf: Mehr Geld und weniger Drogen \newline
    ODER \newline
    Man nimmt das Angebot nicht an/ kauft nichts, man hat also gleich viel Ressourcen und Geld wie davor\\
    \hline
    ID07: Items aufnehmen & Spieler &
    1. Spielfigur befindet sich auf einem Item \newline
    2. Item verschwindet von Map und wird im Inventar abgelegt &
    Spielfigur befindet sich direkt vor dem Item &
    Item befindet sich im Inventar der Spielfigur und ist von der Map verschwunden \\
    \hline
    ID08: Polizist verfolgt Spielfigur & Polizei KI&
    1. Polizist bemerkt verdächtiges Verhalten(z.B. Verkauf von Drogen, schießen, \dots) innerhalb seines Radars \newline
    2. Polizist identifiziert Spieler \newline
    3. Polizist verfolgt Spieler, solange dieser in seinem Sichtfeld/Radar ist  &
    Spielfigur befindet sich im Radar des Polizisten und tut etwas Verdächtiges &
    Spielfigur konnte vor Polizei fliehen oder sich verstecken \newline
    ODER\newline
    Polizei nimmt Spielfigur gefangen und sie kommt ins Gefängnis \\
    \hline
    ID09: verstecken & Spieler &
    1. Spielfigur befindet sich in der Nähe eines Busches\newline
    2. Spieler drückt Taste\newline
    - wenn sich die Spielfigur innerhalb des Radius der Polizei befindet ist verstecken nicht möglich, da die Polizei das sonst sieht
    - außerhalb des Radius der Polizei kann sich die Spielfigur in dem Busch verstecken und ist für die Polizei nicht auffindbar &
    Spielfigur befindet sich neben einem Busch und außerhalb des Radius der Polizei &
    Spielfigur befindet sich in dem Busch und ist für Polizei nicht sichtbar \\
    \hline
    ID10: Waffe wechseln & Spieler &
    1. Spieler drückt Taste zum Waffe wechseln \newline
    2. Waffe wird gewechselt &
    Spielfigur hat Pistole nicht ausgewählt\newline
    ODER\newline
    Spielfigur hat Pistole ausgewählt &

    Spielfigur hat Pistole ausgewählt wenn vorher nicht ausgewählt \newline
    ODER\newline
    Spielfigur hat nicht die Pistole ausgewählt, wenn vorher ausgewählt  \\
    \hline
\end{longtable}

\let\clearpage\origclearpage
\restoregeometry


